\section{Многомерный анализ}

\subsection{Метрика, норма и скалярное произведение}

\begin{defn}{}{}
    Пусть \( X \) --- некоторое множество.

    Функция \( \rho : X \times X \to R_{\geq 0} \) называется метрикой, если:
    \begin{enumerate}
        \item
            \( \rho(x, y) = 0 \Leftrightarrow x = y \)
        \item
            \( \rho(x, y) = \rho(y, x) \)
        \item
            \( \rho(x, y) \leq \rho(x, z) + \rho(z, y) \quad \forall x, y, z \in X \)
    \end{enumerate}

    Пара \( (X, \rho) \) называется метрическим пространством.
\end{defn}

\begin{example}{}{}
    \begin{enumerate}
        \item
            \( X \) --- множество все точек плоскости

            \(
                \rho(x, y)
                =
                \begin{cases}
                    1, x \neq y
                    \\
                    0, x = y
                \end{cases}
            \) --- дискретная метрика
        \item
            \( X \) --- плоскость.

            \( x = \begin{pmatrix} x_1 \\ x_2 \end{pmatrix}, y = \begin{pmatrix} y_1 \\ y_2 \end{pmatrix} \)

            \( \rho_1 (x, y) = |x_1 - y_1| + |x_2 - y_2| \)

            \( \rho_2 (x, y) = \sqrt{(x_1 - y_1)^2 + (x_2 - y_2)^2} \) --- Евклидова метрика
        \item
            \( l^\infty = \{ a_n \} \), причем \( \exists c(a) : \forall n \in \NN \: |a_n| \leq c \)

            \( \rho(a, b) = \sup\limits_{n \in \NN} |a_n - b_n| \)
    \end{enumerate}
\end{example}

\begin{defn}{}{}
    Пусть \( V \) --- векторное пространство над \( \RR \).

    Функция \( || \cdot || : V \to \RR_{\geq 0} \) называется нормой, если:
    \begin{enumerate}
        \item
            \( || \ol{x} || = 0 \Leftrightarrow \ol{x} = \ol{0} \)
        \item
            \( \forall \alpha \in \RR \: || \alpha \ol{x} || = | \alpha | \cdot || \ol{x} || \)
        \item
            \( || \ol{x} + \ol{y} || \leq || \ol{x} || + || \ol{y} || \)
    \end{enumerate}

    Пара \( (V, || \cdot ||) \) --- нормированное пространство.
\end{defn}

\begin{example}{}{}
    \( V = \RR^2 \)

    \begin{itemize}
        \item
            \( || \ol{x} ||_1 = |x_1| + |x_2| \) --- Манхеттенская норма
        \item
            \( || \ol{x} ||_2 = \sqrt{x_1^2 + x_2^2} \) --- Евклидова норма
    \end{itemize}
\end{example}

\begin{defn}{}{}
    Пусть \( V \) --- векторное пространство над \( \RR \).

    Функция \( (\cdot, \cdot) : V \times V \to \RR \) называется скалярным произведением над \( V \), если:
    \begin{enumerate}
        \item
            \( (\ol{x}, \ol{x}) \geq 0 \) и \( (\ol{x}, \ol{x}) = 0 \Leftrightarrow \ol{x} = \ol{0} \)
        \item
            \( (\ol{x}, \ol{y}) = (\ol{y}, \ol{x}) \)
        \item
            \( \alpha, \beta \in \RR \) верно
            \( (\alpha \ol{x} + \beta \ol{y}, \ol{z}) = \alpha (\ol{x}, \ol{z}) + \beta (\ol{y}, \ol{z}) \)
    \end{enumerate}

    Пара \( (V, (\cdot, \cdot)) \) --- Евклидово пространство.
\end{defn}

\begin{thrm}{Неравенство Коши-Буняковского}{}
    \[
        \forall \ol{x}, \ol{y} \in V \quad |(\ol{x}, \ol{y})| \leq \sqrt{(\ol{x}, \ol{x})} \sqrt{(\ol{y}, \ol{y})}
    \]
\end{thrm}

\begin{gather*}
    (\ol{x} - t \ol{y}, \ol{x} - t \ol{y}) \geq 0
    \\
    (\ol{x}, \ol{x}) - 2t (\ol{x}, \ol{y}) + t^2 (\ol{y}, \ol{y}) \geq 0
\end{gather*}

\begin{enumerate}
    \item
        \( (\ol{y}, \ol{y}) = 0 \)

        Тогда \( \ol{y} = 0 \).

        Покажем, что \( (\ol{x}, \ol{y}) = 0 \).

        \( (\ol{x}, \ol{0}) = (\ol{x}, \ol{0} + \ol{0}) = 2 (\ol{x}, \ol{0}) \).

        Значит скалярное произведение с \( \ol{0} \) --- нулевое,
        неравенство верно.
    \item
        \( (\ol{y}, \ol{y}) \neq 0 \)

        Тогда \( D \leq 0 \)
        \[
            D = 4 (\ol{x}, \ol{y})^2 - 4(\ol{x}, \ol{x})(\ol{y}, \ol{y}) \leq 0
        \]

        А это и есть нужное неравенство.
\end{enumerate}

\newpage

\begin{thrm}{}{}
    \( || \ol{x} || = \sqrt{(\ol{x}, \ol{x})} \) задает норму \( V \)
\end{thrm}

Остается проверить только третье свойство нормы (остальные два очевидны):
\begin{gather*}
    || \ol{x} + \ol{y} ||^2 = (\ol{x} + \ol{y}, \ol{x} + \ol{y}) = (\ol{x}, \ol{x}) + 2(\ol{x}, \ol{y}) + (\ol{y}, \ol{y})
    \leq
    \\
    \leq
    (\ol{x}, \ol{x}) + 2 \sqrt{(\ol{x}, \ol{x})} \sqrt{(\ol{y}, \ol{y})} + (\ol{y}, \ol{y})
    =
    \left( \sqrt{(\ol{x}, \ol{x})} + \sqrt{(\ol{y}, \ol{y})} \right)^2
    =
    \\
    =
    \left( || \ol{x} || + || \ol{y} || \right)^2
\end{gather*}

\begin{defn}{}{}
    \[
        \cos \varphi = \frac{(\ol{x}, \ol{y})}{||\ol{x}|| \cdot ||\ol{y}||}
    \]

    При этом считаем, что нулевой вектор ортогонален всем остальным.
\end{defn}

\begin{example}{}{}
    Пусть \( V = \RR^d \)

    \(
        \ol{x} = \begin{pmatrix}
            x_1
            \\
            \vdots
            \\
            x_d
        \end{pmatrix}
    \),
    \(
        \ol{y} = \begin{pmatrix}
            y_1
            \\
            \vdots
            \\
            y_d
        \end{pmatrix}
    \)

    \begin{itemize}
        \item
            \(
                \displaystyle
                (\ol{x}, \ol{y}) = \sum\limits_{j = 1}^{d} x_j y_j
            \)

            Тогда
            \(
                \displaystyle
                \sqrt{(\ol{x}, \ol{x})} = \sqrt{\sum\limits_{j = 1}^{d} x_j^2}
            \)
        \item
            \(
                \displaystyle
                || \ol{x} ||_1 = \sum\limits_{j = 1}^d |x_j|
            \) --- не порождается скалярным произведением
    \end{itemize}
\end{example}

\begin{thrm}{(Йордан -- Фон Нейман)}{}
    \begin{gather*}
        \text{\( || \cdot || \) на \( V \) задается скалярным произведением}
        \\
        \Updownarrow
        \\
        \forall \ol{x}, \ol{y} \in V \quad 2 ||\ol{x}||^2 + 2 ||\ol{y}||^2 = || \ol{x} + \ol{y} ||^2 + || \ol{x} - \ol{y} ||^2
    \end{gather*}
\end{thrm}

\subsection{Предел, непрерывность и связность}

\begin{defn}{}{}
    \begin{enumerate}
        \item
            Множество \( B_r(a) = \{ x \in X : \rho(x, a) < r \} \) называется открытым шаром в \( X \)
            радиуса \( r \) с центром в точке \( a \).
        \item
            Множество \( \ol{B}_r(a) = \{ x \in X : \rho(x, a) \leq r \} \) называется замкнутым шаром в \( X \)
            радиуса \( r \) с центром в точке \( a \).
        \item
            \( B_{\varepsilon}'(a) = \mathring{B_{\varepsilon}} (a) = \{ x \in X : 0 < \rho(x, a) < \varepsilon \} \)
            --- проколотая окрестность точки \( a \)
    \end{enumerate}
\end{defn}

Научимся делать так, чтобы шар с большим радиусом был внутри шара с меньшим радиусом.
Возьмем три точки \( A, B, C \), образующие прямоугольный треугольник:
\begin{itemize}
    \item
        \( AB = 3 \)
    \item
        \( BC = 4 \)
    \item
        \( AC = 5 \)
\end{itemize}

Возьмем \( B_{4.1}(B) = \{ A, B, C \} \) и \( B_{4.9}(C) = \{ B, C \} \) --- ура.

\begin{exercise}{}{}
    Доказать, что в нормированном пространстве шар с большим радиусом не может содержаться в шаре с меньшим радиусом.
\end{exercise}

\begin{defn}{}{}
    Пусть \( (X, \rho) \) --- метрическое пространство.

    Функция \( f : \NN \to X \) называется последовательностью элементов из \( X \).

    Будем обозначать \( f(n) = x_n \in X \).

    \( \lim x_n = x \Leftrightarrow \exists \varepsilon > 0 \ \exists N : \forall n > N \: \rho(x_n, x) < \varepsilon \)
\end{defn}

\begin{defn}{}{}
    Пусть множество \( M \subset X \)

    \begin{itemize}
        \item
            \( a \) --- предельная точка \( M \), если \( \forall \varepsilon > 0 \ M \cap \mathring{B}_{\varepsilon}(a) \neq \varnothing \)
        \item
            \( a \in M \) --- изолированная точка множества \( M \), если \( a \) не является предельной.
        \item
            \( a \) --- граничная точка, если в любой ее окрестности есть точки и из \( M \), и из \( X \setminus M \)
    \end{itemize}
\end{defn}

\begin{defn}{}{}
    \begin{itemize}
        \item
            Множество \( B \subset X \) открыто, если \( \forall a \in B \ \exists \varepsilon > 0 : B_{\varepsilon} (a) \subset B \)
        \item
            Множество \( C \subset X \) замкнуто, если \( X \setminus C \) открыто.
    \end{itemize}
\end{defn}

\begin{thrm}{}{}
    Множество \( C \subset X \) замкнуто тогда и только тогда, когда оно содержит все свои предельные точки.
\end{thrm}

Доказательство такое же, как и в первом семестре.

\begin{thrm}{}{}
    Пусть \( (X, \rho) \) --- метрическое пространство и \( x_n \to x, y_n \to y \) в \( X \)

    Тогда:
    \begin{enumerate}
        \item
            \( \lim\limits_{n \to \infty} \rho(x_n, y_n) = \rho(x, y) \)
        \item
            \( \exists! \lim\limits_{n \to \infty} x_n \)
        \item
            \( \{ x_n \} \) ограничено, то есть содержится в некотором шаре.
    \end{enumerate}
\end{thrm}

\begin{enumerate}
    \item
        \begin{gather*}
            | \rho(x_n, y_n) - \rho(x, y) |
            \leq
            | \rho(x_n, y_n) - \rho(x_n, y) | + | \rho(x_n, y) - \rho(x, y) |
            \leq
            \\
            \leq
            \rho(y_n, y) + \rho(x_n, x)
            <
            2 \varepsilon
        \end{gather*}
    \item
        Если есть два предела, то по предыдущему между ними нулевое расстояние,
        то есть они равны.
    \item
        Доказательство такое же, как и с обычными пределами (суффикс ограничиваем по определению предела,
        расширяем шар так, чтобы префикс тоже в него попал).
\end{enumerate}

\begin{thrm}{}{}
    Пусть \( (V, || \cdot ||) \) --- нормированное пространство.

    Пусть \( \ol{x}_n \to \ol{x}, \ol{y}_n \to \ol{y} \).

    Пусть \( \{ \alpha_n \} \subset \RR : \alpha_n \to \alpha \in \RR \).

    Тогда
    \begin{enumerate}
        \item
            \( \lim\limits_{n \to \infty} (\ol{x}_n + \ol{y}_n) = \ol{x} + \ol{y} \)
        \item
            \( \lim\limits_{n \to \infty} \alpha_n \ol{x}_n = \alpha \ol{x} \)
    \end{enumerate}
\end{thrm}

\begin{enumerate}
    \item
        Тривиально
    \item
        \(
            \forall \varepsilon > 0 \ \exists N : \forall n > N \
            |\alpha_n - \alpha| < \varepsilon, || \ol{x}_n - \ol{x} || < \varepsilon
        \)

        \begin{gather*}
            || \alpha_n \ol{x}_n - \alpha \ol{x} ||
            \leq
            || \alpha_n \ol{x}_n - \alpha_n \ol{x} || + || \alpha_n \ol{x} - \alpha \ol{x} ||
            =
            \\
            =
            | \alpha_n | \cdot || \ol{x}_n - \ol{x} || + |\alpha_n - \alpha| \cdot || \ol{x} ||
            \leq
            (| \alpha_n | + || \ol{x} ||) \cdot \varepsilon
        \end{gather*}

        Заметим, что \( |\alpha_n| \) ограничено в силу сходимости, а \( || \ol{x} || \) --- константа.
        А значит существует \( C > 0 : | \alpha_n | + || \ol{x} || < C \):
        \[
            || \alpha_n \ol{x}_n - \alpha \ol{x} ||
            \leq
            (| \alpha_n | + || \ol{x} ||) \cdot \varepsilon
            <
            C \cdot \varepsilon
        \]

        Сделав \( \varepsilon \) достаточно маленьким, докажем нужное неравенство.
\end{enumerate}

\begin{thrm}{}{}
    \( V \) --- Евклидово пространство, \( || \cdot || \) порождена скалярным произведением.

    Пусть \( \lim\limits_{n \to \infty} \ol{x}_n = \ol{x} \) и \( \lim\limits_{n \to \infty} \ol{y}_n = \ol{y} \)

    Тогда \( \lim\limits_{n \to \infty} (\ol{x}_n, \ol{y}_n) = (\ol{x}, \ol{y}) \)
\end{thrm}

\begin{gather*}
    | (\ol{x}_n, \ol{y}_n) - (\ol{x}, \ol{y}) |
    \leq
    | (\ol{x}_n, \ol{y}_n) - (\ol{x}, \ol{y}_n)| + | (\ol{x}, \ol{y}_n) - (\ol{x}, \ol{y}) |
    =
    \\
    =
    | (\ol{x}_n - \ol{x}, \ol{y}_n) | + | (\ol{x}, \ol{y}_n - \ol{y}) |
    \leq
    || \ol{x}_n - \ol{x} || \cdot || \ol{y}_n || + || \ol{x} || \cdot || \ol{y}_n - \ol{y} ||
\end{gather*}

\( \forall \varepsilon > 0 \: \exists N : \forall n > N \: || \ol{x}_n - \ol{x} || < \varepsilon \)
и \( || \ol{y}_n - \ol{y} || < \varepsilon \)

\( || \ol{x} || \) --- константа, а \( || \ol{y}_n || \) ограничено, так как имеет предел.
Значит их сумма тоже ограничена некоторой константой \( C > 0 \).

\begin{gather*}
    | (\ol{x}_n, \ol{y}_n) - (\ol{x}, \ol{y}) |
    \leq
    || \ol{x}_n - \ol{x} || \cdot || \ol{y}_n || + || \ol{x} || \cdot || \ol{y}_n - \ol{y} ||
    <
    \\
    <
    ( || \ol{y}_n || + || \ol{x} || ) \cdot \varepsilon
    <
    C \cdot \varepsilon
\end{gather*}

Что и требовалось доказать.

\begin{thrm}{}{}
    Пусть \( \ol{x_n} \to x \) в \( V = \RR^d \) (в евклидовой норме)

    Тогда \( x_j^n \to x_j \) (\( x_j^n \) --- \( j \)-я координата \( \ol{x_n} \))

    (верно в обе стороны)
\end{thrm}

\[
    || \ol{x_n} - \ol{x} || = \sqrt{\sum\limits_{j = 1}^{d} (x_j^n - x_j)^2}
\]

Пусть \( C_n = \max\limits_{j = 1, \ldots, d} (x_j^n - x_j)^2 \)

\[
    C_n \leq || \ol{x_n} - \ol{x} || \leq \sqrt{d} \cdot \sqrt{C_n}
\]

Значит из поокоординатной сходимости следует просто сходимость, и наоборот.

\begin{thrm}{(Больцано-Вейерштрасса)}{}
    Пусть \( \{ \ol{x_n} \} \subset \RR^d \)

    Пусть \( \{ \ol{x_n} \} \) ограничен \( \Rightarrow \) есть сходящаяся подпоследовательность.

    Считаем, что у нас Евклидова норма, но позже поймем, что это неважно.
\end{thrm}

Очевидно, что последовательность каждой из координат ограничена.
Применим ``одномерную'' теорему Больцано-Вейерштрасса для первой координаты.
Так мы выделим подпоследовательность, у которой сходится первая координата.
Теперь на этой подпоследовательности вновь применим данную теорему,
получим подпоследовательность изначальной последовательности,
у которой сходится уже две первых координаты.
Сделав так для всех оставшихся координат, докажем теорему.

\begin{defn}{(предел по Коши)}{}
    Пусть \( (X, \rho_x) \) и \( (Y, \rho_y) \) --- метрические пространства.

    \( a \) --- предельная точка \( X \).

    \( f : X \setminus \{ a \} \to Y \)

    Приведем две равносильных формулировки:
    \begin{itemize}
        \item
            \(
                \lim\limits_{x \to a} f(x) = A
                \Leftrightarrow
                \forall \varepsilon > 0 \: \exists \delta > 0 : \forall x \in X : 0 < \rho_x(a, x) < \delta \: \rho_y(A, f(x)) < \varepsilon
            \)
        \item
            \(
                \lim\limits_{x \to a} f(x) = A
                \Leftrightarrow
                \forall \varepsilon > 0 \: \exists \delta > 0 : \forall x \in \mathring{B}_{\delta}(a) \: f(x) \in B_{\varepsilon}(A)
            \)
    \end{itemize}
\end{defn}

\begin{defn}{(предел по Гейне)}{}
    Пусть \( (X, \rho_x) \) и \( (Y, \rho_y) \) --- метрические пространства.

    \( a \) --- предельная точка \( X \).

    \( f : X \setminus \{ a \} \to Y \)

    \(
        \lim\limits_{x \to a} f(x) = A
        \Leftrightarrow
        \forall \{ a_n \} \subset X : \forall n \in \NN \: a_n \neq a, \lim\limits_{n \to \infty} a_n = a
        \text{ верно } \lim\limits_{n \to \infty} f(a_n) = A
    \)
\end{defn}

\begin{thrm}{}{}
    Определения по Коши и Гейне равносильны.
\end{thrm}

\newpage

\begin{enumerate}
    \item
        \( \Rightarrow \)

        Пусть \( \{ a_n \} \subset X : \forall n \in \NN \: a_n \neq a, \lim\limits_{n \to \infty} a_n = a \).

        То есть
        \(
            \forall \delta > 0 \: \exists N > 0 : \forall n > N \: 0 < \rho_x (a, a_n) < \delta
            \Rightarrow
            \rho_y (A, f(a_n)) < \varepsilon
        \)
    \item
        \( \Leftarrow \)

        Пусть определение по Коши не выполнено.

        \( \exists \varepsilon > 0 : \forall \delta > 0 \exists x \in \mathring{B}_{\delta} : \rho_y(f(x), A) \geq \varepsilon \)

        Выберем при \( \delta = 1 \) \( x_1 : \rho_y (f(x_1), A) \geq \varepsilon \)

        При \( \delta = \frac{1}{2} \) возьмем \( x_2 : \rho_y (f(x_2), A) \geq \varepsilon \)

        И так далее

        Но тогда для \( \{ x_n \} \) неверно определение предела по Гейне --- противоречие.
\end{enumerate}

\begin{thrm}{}{}
    Пусть \( f, g : (X, \rho_x) \to (V, || \cdot ||) \) и \( \alpha : (X, \rho_x) \to \RR \)

    \begin{enumerate}
        \item
            \( \lim\limits_{x \to a} f(x) = A \)
        \item
            \( \lim\limits_{x \to a} g(x) = B \)
        \item
            \( \lim\limits_{x \to a} g(x) = \alpha_0 \)
    \end{enumerate}

    Тогда:
    \begin{enumerate}
        \item
            \( \lim\limits_{x \to a} (f(x) + g(x)) = A + B \)
        \item
            \( \lim\limits_{x \to a} \alpha(x) \cdot f(x) = \alpha_0 \cdot A \)
    \end{enumerate}
\end{thrm}


\begin{thrm}{}{}
    Пусть \( f, g : (X, \rho_x) \to (V, (\cdot, \cdot) ) \)

    \begin{enumerate}
        \item
            \( \lim\limits_{x \to a} f(x) = A \)
        \item
            \( \lim\limits_{x \to a} g(x) = B \)
    \end{enumerate}

    Тогда \( \lim\limits_{x \to a} (f(x), g(x)) = (A, B) \)
\end{thrm}

Мгновенно доказывается через определение по Гейне.

\begin{thrm}{}{}
    Пусть \( f : (X, \rho_x) \to (Y, \rho_y) \) и \( \lim\limits_{x \to a} f(x) = A \)

    \begin{enumerate}
        \item
            \( \exists \delta, r > 0 : \forall x \in \mathring{B}_\delta(a) \: f(x) \in B_r(A) \).
        \item
            \( \forall B \neq A \: \exists \delta > 0 : \forall x \in \mathring{B}_\delta(a) \: f(x) \notin B_{\frac{r}{2}}(B) \),
            где \( r = \rho_y(B, A) \)
    \end{enumerate}
\end{thrm}

\begin{enumerate}
    \item
        Явно следует из определение предела по Коши.
    \item
        Возьмем \( \varepsilon = \frac{r}{2} \).

        Из определения предела по Коши существует \( \delta : \forall x \in \mathring{B}_{\delta} \ f(x) \in B_{\frac{r}{2}} (A) \).

        Понятно, что \( B_{\frac{r}{2}} (A) \cap B_{\frac{r}{2}} (B) = \varnothing \).
\end{enumerate}

\begin{defn}{Непрерывность в точке по Гейне}{}
    Пусть \( f : (X, \rho_x) \to (Y, \rho_y) \), \( a \in X \)

    \( f \in C(a) \Leftrightarrow \forall \{ a_n \} \subset X : \lim\limits_{n \to \infty} a_n = a \)
    выполнено \( \lim\limits_{n \to \infty} f(a_n) = f(a) \)
\end{defn}

\begin{defn}{Непрерывность композиции}{}
    Пусть \( f : (X, \rho_x) \to (Y, \rho_y) \), \( g : (Y, \rho_y) \to (Z, \rho_z) \)

    Также \( f \in C(a), g \in C(f(a)) \).

    Тогда \( (g \circ f) \in C(a) \).
\end{defn}

\( \forall \{ a_n \} \subset X : a_n \to a \lim\limits_{n \to \infty} f(a_n) = f(a) \)

Тогда \( \lim\limits_{n \to \infty} g(f(a_n)) = g(f(a)) \)

Что и требовалось доказать.

\begin{defn}{Непрерывность в точке по Коши}{}
    Пусть \( f : (X, \rho_x) \to (Y, \rho_y) \), \( a \in X \)

    \( f \in C(a) \Leftrightarrow \forall \varepsilon > 0 \: \exists \delta > 0 : \forall x \in B_{\delta} (x) \: f(x) \in B_{\varepsilon} (f(a)) \)
\end{defn}

\begin{defn}{Непрерывность на множестве}{}
    Пусть \( f : (X, \rho_x) \to (Y, \rho_y) \), \( M \subset X \)

    \( f \) непрерывна на \( M \),
    если она непрерывна в каждой точке этого множества.
\end{defn}

\newpage

\begin{defn}{Компакт}{}
    Множество \( K \subset X \) называется компактом тогда и только тогда,
    когда из любого его покрытия открытыми множествами можно выбрать конечное подпокрытие.
\end{defn}

\begin{thrm}{Критерий компактности в \( R^d \)}{}
    Множество \( K \subset R^d \) --- компакт тогда и только тогда,
    когда оно замкнуто и ограничено.
\end{thrm}

\begin{thrm}{}{}
    Пусть \( K \subset \RR^d \) --- компакт, \( f : \RR^d \to \RR^m, f \in C(K) \),
    тогда \( f(K) \) --- тоже компакт (верно и для произвольных метрических пространств).
\end{thrm}

\begin{defn}{Равномерная непрерывность}{}
    \( f : (X, \rho_x) \to (Y, \rho_y) \)

    \( f \) равномерно непрерывна на множестве \( M \) тогда и только тогда,
    когда \( \forall \varepsilon > 0 \: \exists \delta > 0 : \forall x, y \in M : \rho_x(x, y) < \delta \)
    выполнено \( \rho_y (f(x), f(y)) < \varepsilon \)
\end{defn}

\begin{thrm}{(Гейне-Кантора)}{}
    Функция, непрерывная на компакте \( K \subset R^d \), равномерно непрерывна на нем
    (верно для произвольных метрических пространств).
\end{thrm}

\begin{defn}{}{}
    Множество \( E \subset V \) называется связным,
    если не существует открытых множеств \( A, B \subset V \) таких,
    что \( A \cap M \neq \varnothing, B \cap M \neq \varnothing \) и \( M \subset A \sqcup B \).
\end{defn}

\begin{thrm}{(Больцано-Коши)}{}
    Пусть \( f : \RR^d \to \RR^m \), \( f \in C(M) \), причем \( M \) связно.
    Тогда \( f(M) \) тоже связно.
\end{thrm}

\begin{defn}{}{}
    Множество \( D \subseteq V \) линейно связно,
    если \( \forall a, b \in D \) существует непрерывная функция \( \alpha : [0, 1] \to D \)
    такая, что \( \alpha(0) = a, \alpha(1) = b \)
\end{defn}

\begin{exercise}{}{}
    \begin{enumerate}
        \item
            Если множество \( D \subseteq V \) линейно связно, то \( D \) связно
        \item
            Если \( E \subseteq V \) открыто и связно, то \( E \) линейно связно
        \item
            Пусть \( E \subseteq \RR^2 \) задается так:
            \[
                \left\{ \left( x, \sin \frac{1}{x} \right), x \neq 0 \right\} \cup \{ (x, y) : x = 0, y \in [-1, 1] \}
            \]

            \( E \) связно, но не линейно связно
    \end{enumerate}
\end{exercise}

\begin{defn}{}{}

    Пусть \( V \) --- линейное пространство и \( || \cdot ||_1 \), \( || \cdot ||_2 \) заданы на \( V \)

    Нормы называются эквивалентными, если \( \exists c_1, c_2 > 0 \) такие,
    что \( \forall \ol{x} \in V \) выполнено
    \(
        c_1 || \ol{x} ||_1 \leq || \ol{x} ||_2 \leq c_2 || \ol{x} ||_1
    \)
\end{defn}

\begin{thrm}{}{}
    Все нормы на конечномерном пространстве \( V \) эквивалентны.
\end{thrm}

Понятно, что определение действительно дает нам отношение эквивалентности.
Будем доказывать эквивалентность всех норм манхеттенской.

Пусть \( V = \langle \ol{e}_1, \ldots, \ol{e}_d \rangle \),
где \( \ol{e}_1, \ldots, \ol{e}_d \) --- линейно независимы.

Пусть \( \ol{x} = \sum\limits_{j = 1}^d x_j \ol{e}_j \)

Пусть \( p(\ol{x}) = |x_1| + \ldots + |x_d| \)

Пусть \( q \) --- другая норма на \( V \), и пусть \( c_2 = \max\limits_{j = 1, \ldots, d} q(\ol{e}_j) > 0 \).

Тогда \( q(\ol{x}) \leq \sum\limits_{j = 1}^{d} |x_j| \cdot q(\ol{e}_j) \leq c_2 \cdot p(\ol{x}) \)

По неравенству треугольника:
\( | q(\ol{x}) - q(\ol{y}) | \leq q(\ol{x} - \ol{y}) \leq c_2 \cdot p(\ol{x} - \ol{y}) \),
поэтому \( q : (V, p) \to \RR \) непрерывна на \( V \),
а значит непрерывна и на сфере \( S \) радиуса \( 1 \).
Но \( S \) --- компакт, а значит \( \exists c_1 = \min\limits_{\ol{x} \in S} q(\ol{x}) > 0 \).

Значит \( \forall \ol{x} \neq \ol{0} \) верно
\(
    q \left( \frac{\ol{x}}{p(\ol{x})} \right) \geq c_1 = c_1 p \left( \frac{\ol{x}}{p(\ol{x})} \right)
    \Rightarrow
    q ( \ol{x} ) \geq c_1 p ( \ol{x} )
\)

То есть \( \forall \ol{x} \neq \ol{0} \) верно \( c_1 p(\ol{x}) \leq q(\ol{x}) \leq c_2 p(\ol{x}) \),
а для \( \ol{x} = 0 \) это неравенство, очевидно, верно \( \Rightarrow \)
любая норма эквивалентна \( p \), а значит все нормы эквивалентны.

\begin{exercise}{}{}
    В бесконечномерном пространстве найдутся две неэквивалентные нормы.

    (нужно использовать базис Гамеля)
\end{exercise}

\newpage

\subsection{Дифференциал}

\( f : \RR^d \to \RR^m \)

\(
    f(\ol{x}) = \begin{pmatrix}
        f_1 (\ol{x})
        \\
        \vdots
        \\
        f_m (\ol{x})
    \end{pmatrix}
\)

\(
    f_j : \RR^d \to \RR
\) --- координатная функция

\begin{defn}{}{}
    Пусть \( f : \RR^d \to \RR^m \)

    \( f \in D(\ol{a}) \),
    если \( \exists \mathring{U}(\ol{0}) \subset \RR^d \),
    \( \exists L_{\ol{a}} : \RR^d \to \RR^m \) такие,
    что
    \begin{gather*}
        \forall h \in \mathring{U}(\ol{0}) \
            f(\ol{a} + \ol{h}) - f(\ol{a}) = L_{\ol{a}} \ol{h} + || \ol{h} || \cdot \alpha(\ol{h})
        \\
        \alpha : \mathring{V}(\ol{0}) \to \RR^m \quad \lim\limits_{\ol{h} \to \ol{0}} \alpha(\ol{h}) = \ol{0} \in \RR^m
    \end{gather*}

    \( L_{\ol{a}} \) называется дифференциалом \( f \) в точке \( \ol{a} \) и обозначается \( df(\ol{a}) \)
\end{defn}

\begin{example}{}{}
    \[
        f \begin{pmatrix}
            x_1
            \\
            x_2
        \end{pmatrix}
        =
        \begin{pmatrix}
            x_1^2
            \\
            x_2^3
        \end{pmatrix}
        \quad
        f : \RR^2 \to \RR^2
    \]

    Подставим:
    \begin{gather*}
        f \begin{pmatrix}
            \ol{a} + \ol{h}
        \end{pmatrix}
        =
        \begin{pmatrix}
            (a_1 + h_1)^2
            \\
            (a_2 + h_2)^3
        \end{pmatrix}
        =
        \begin{pmatrix}
            a_1^2 + 2 a_1 h_1 + h_1^2
            \\
            a_2^3 + 3 a_2^2 h_2 + 3 a_2 h_2^2 + h_2^3
        \end{pmatrix}
        =
        \\
        =
        \begin{pmatrix}
            a_1^2
            \\
            a_2^3
        \end{pmatrix}
        +
        \begin{pmatrix}
            2a_1 & 0
            \\
            0 & 3a_2^2
        \end{pmatrix}
        \begin{pmatrix}
            h_1
            \\
            h_2
        \end{pmatrix}
        +
        \sqrt{h_1^2 + h_2^2}
        \begin{pmatrix}
            \frac{h_1^2}{\sqrt{h_1^2 + h_2^2}}
            \\
            \frac{3 a_2 h_2^2 + h_2^3}{\sqrt{h_1^2 + h_2^2}}
        \end{pmatrix}
    \end{gather*}
\end{example}

\begin{prop}{}{}
    \( f \in D(\ol{a}) \Rightarrow f \in C(\ol{a}) \)
\end{prop}

\begin{gather*}
    ||f( \ol{a} + \ol{h} ) - f(\ol{a})||
    =
    || L_{\ol{a}} \ol{h} + || \ol{h} || \cdot \alpha(\ol{h}) ||
    \leq
    \\
    \leq
    \sum\limits_{j = 1}^{d} |h_j| \cdot || L_{\ol{a}} \ol{e}_j ||
        + || \ol{h} || \cdot || \alpha(\ol{h}) ||
    \to \ol{0} \quad (\ol{h} \to \ol{0})
\end{gather*}

\begin{prop}{}{}
    Пусть \( f \in D(\ol{a}) \)

    Тогда \( \exists! \: df(\ol{a}) \)
\end{prop}

Пусть \( \exists L_1, L_2 \) такие, что
\begin{enumerate}
    \item
        \( f(\ol{a} + \ol{h}) - f(\ol{a}) = L_1 \ol{h} + ||\ol{h} || \alpha_1 (\ol{h}) \)
    \item
        \( f(\ol{a} + \ol{h}) - f(\ol{a}) = L_2 \ol{h} + ||\ol{h} || \alpha_2 (\ol{h}) \)
\end{enumerate}

Пусть \( L = L_2 - L_1 \).
\begin{gather*}
    L \ol{h} = ( \alpha_1(\ol{h}) -\alpha_2(\ol{h}) ) || \ol{h} ||
    \\
    L \frac{\ol{h}}{||h||} = \alpha_1(\ol{h}) -\alpha_2(\ol{h})
\end{gather*}

Если взять предел \( \ol{h} \to \ol{0} \),
то величина справа стремится к нулю.
Но тогда \( L \) равно \( \ol{0} \) на любом векторе единичной длины,
а значит \( L = 0 \).

\begin{prop}{}{}
    \[
        f \in D(\ol{a})
        \Leftrightarrow
        \exists L_{\ol{a}} : \ \lim\limits_{|| \ol{h} || \to 0} \frac{||f(\ol{a} + \ol{h}) - f(\ol{a}) - L_{\ol{a}} \ol{h}||}{||\ol{h}||} = 0
    \]
\end{prop}

\begin{itemize}
    \item
        \( \Rightarrow \)

        Следует из определения дифференцируемости в точке
    \item
        \( \Leftarrow \)

        \( \frac{f(\ol{a} + \ol{h}) - f(\ol{a}) - L_{\ol{a}} \ol{h}}{||\ol{h}||} = \alpha(\ol{h}) \)

        Тогда \( \lim\limits_{|| \ol{h} || \to 0} \alpha(\ol{h}) = \ol{0} \).

        Домножим на \( || \ol{h} || \), получим искомое
\end{itemize}

\begin{remark}{}{}
    Строчки матрицы дифференциалов являются дифференциалами координатных функций.
\end{remark}

\newpage

\subsection{Производная по направлению}

\begin{defn}{Производная по направлению (вектору)}{}
    Если существует предел \( \lim\limits_{t \to 0} \frac{f(\ol{a} + t \cdot \ol{v}) - f(\ol{a})}{t} \),
    то он называется производной \( f \) по вектору \( \ol{v} \).
    Если \( ||\ol{v}|| = 1 \), то этот предел называется производной по направлению \( \ol{v} \)
    и обозначается \( f_t ' (\ol{a} + t \cdot \ol{v}) \: \vline_{\: t = 0} \)
    или \( \frac{\partial f}{\partial \ol{v}} (\ol{a}) \)
\end{defn}

\begin{thrm}{}{}
    \[
        f \in D(\ol{a}) \Rightarrow \lim\limits_{t \to 0} \frac{f(\ol{a} + t \ol{v}) - f(\ol{a})}{t} = df(\ol{a})(\ol{v})
    \]
\end{thrm}

\begin{gather*}
    f(\ol{a} + t \ol{v}) - f(\ol{a}) = df(\ol{a}) (t \ol{v}) + \alpha(t \ol{v}) \cdot || t \ol{v} ||
    \\
    \frac{f(\ol{a} + t \ol{v}) - f(\ol{a})}{t} = df(\ol{a}) (\ol{v}) + sign(t) \cdot \alpha(t \ol{v}) \cdot || \ol{v} ||
\end{gather*}

\subsection{Частная производная и запись дифференциала с их помощью}

\begin{defn}{Частная производная}{}
    \( f : \RR^d \to \RR \)

    Если \( \exists \lim\limits_{t \to 0} \frac{f(\ol{a} + t e_j)  - f(\ol{a})}{t} \),
    где \( e_j \) --- элемент стандартного базиса, то его называют частной производной \( f \) в точке \( \ol{a} \)
    по переменной \( x_j \).

    Обозначается \( f_{x_j}' (\ol{a}) \) или \( \displaystyle \frac{\partial f}{\partial x_j} (\ol{a}) \)
\end{defn}

\begin{example}{}{}
    \( f(x_1, x_2, x_3) = e^{x_1 x_2} \sin x_3 \).

    Найдем все частные производные:
    \begin{itemize}
        \item
            \( f_{x_1}' (x_1, x_2, x_3) = x_1 e^{x_1 x_2} \sin x_3 \)
        \item
            \( f_{x_2}' (x_1, x_2, x_3) = x_2 e^{x_1 x_2} \sin x_3 \)
        \item
            \( f_{x_3}' (x_1, x_2, x_3) = e^{x_1 x_2} \cos x_3 \)
    \end{itemize}
\end{example}

\begin{example}{}{}
    \[
        f(x, y) = \begin{cases}
            0, \quad x = 0 \lor y = 0
            \\
            1, \quad \text{иначе}
        \end{cases}
    \]

    У этой функции есть обе частные производные в нуле,
    однако ни по какому другому направлению производной не существует.
\end{example}

В случае функции \( f : \RR^d \to \RR \),
\( df(\ol{a}) = (A_1, \ldots, A_d) \).
Поэтому по определению дифференциала:
\[
    f(\ol{a} + \ol{h}) - f(\ol{a}) = A_1 h_1 + \ldots + A_d h_d + o(\ol{h})
\]

Если подставить \( h \) такое, что все его координаты кроме \( h_i \) равны \( 0 \),
то получим
\(
    f(a_1, \ldots, a_{i - 1}, a_i + h_i, a_{i + 1}, \ldots a_d) - f(a_1, \ldots, a_{i - 1}, a_i, a_{i + 1}, \ldots a_d)
    =
    A_i h_i + o(\ol{h})
\).
Откуда \( A_i = \frac{\partial f}{\partial x_i} (\ol{a}) \).

А значит \( df(\ol{a}) = \left( \frac{\partial f}{\partial x_1} (\ol{a}), \ldots, \frac{\partial f}{\partial x_d} (\ol{a}) \right) \)

\begin{remark}{}{}
    \(
        \displaystyle
        \ol{h}_j (\ol{e}_j) = \delta_{ji} = \begin{cases}
            1, j = i
            \\
            0, j \neq i
        \end{cases}
    \)

    \(
        d \ol{h}_j (\ol{x}) (\ol{h}) = x_j + h_j - x_j = h_j
    \)

    Такие дифференциалы обозначаются \( dx_j \)

    То есть \( \forall h \in \RR^d \ dx_j (\ol{h}) = h_j \)
\end{remark}

\begin{remark}{}{}
    \[
        df (\ol{a}) \ol{h}
        =
        \sum\limits_{j = 1}^d \frac{\partial f}{\partial x_j} (\ol{a}) (dx_j (\ol{h}))
        =
        \sum\limits_{j = 1}^d \frac{\partial f}{\partial x_j} (\ol{a}) h_j
    \]
\end{remark}

\begin{example}{}{}
    \( f(x, y) = \sqrt{ |xy| } \)

    \( f_x'(0, 0) = \lim\limits_{h_1 \to 0} \frac{\sqrt{h_1 \cdot 0} - 0}{h_1} = 0 \)

    \( f_y'(0, 0) = \lim\limits_{h_2 \to 0} \frac{\sqrt{0 \cdot h_2} - 0}{h_2} = 0 \)

    \( f(\ol{0} + \ol{h}) - f(\ol{0}) = df(\ol{0}) \ol{h} + o(\ol{h}) = o(\ol{h}) \)

    Пусть
    \(
        \ol{h} = \begin{pmatrix}
            h
            \\
            h
        \end{pmatrix},
        h > 0
        \Rightarrow
        f(\ol{0} + \ol{h}) - f(\ol{0}) = h \neq o(\ol{h})
    \) --- противоречие.
\end{example}

\begin{defn}{Градиент}{}
    \[
        \nabla f (\ol{a}) = \left( \frac{\partial f}{\partial x_1} (\ol{a}), \ldots, \frac{\partial f}{\partial x_d} (\ol{a}) \right)
    \]
\end{defn}

\begin{thrm}{}{}
    Пусть \( f \in D(\ol{a}) \).

    Тогда максимальное значение производной по направлению \( \frac{\partial f}{\partial \ol{l}} (\ol{a}) \)
    будет достигаться при \( \ol{l} = \frac{\nabla f(\ol{a})}{|| \nabla f(\ol{a}) ||} \) (\( \nabla f(\ol{a}) \neq 0 \))
\end{thrm}

\[
    \frac{\partial f}{\partial \ol{l}} (\ol{a}) = df(\ol{a}) \ol{l} = (\nabla f(\ol{a}), \ol{l})
    \leq
    || \nabla f(\ol{a}) || \cdot || \ol{l} ||
    =
    || \nabla f(\ol{a}) ||
\]

А равенство здесь по КБШ достигается как раз при сонаправленности.
Поскольку \( || \ol{l} || = 1 \), доказали теорему.

\begin{defn}{}{}
    \( f : \RR^2 \to \RR \)

    \( \{ (x, y) \: \vline \: f(x, y) = c \} \) --- множество \( c \)-уровня.
\end{defn}

Из теоремы понятна идея градиентного спуска: нужно на ``немного'' сдвигать точку
по вектору, противоположному градиенту (однако можем попасть в локальный минимум).

\begin{defn}{Матрица Якоби}{}
    Пусть существуют \( \frac{\partial f_j}{\partial x_i} (\ol{a}) \),
    где \( j = 1, \ldots, m \), а \( i = 1, \ldots, d \).
    \[
        J_f (\ol{a}) = \begin{pmatrix}
            \frac{\partial f_1}{\partial x_1} (\ol{a}) & \ldots & \frac{\partial f_1}{\partial x_d} (\ol{a})
            \\
            \vdots & & \vdots
            \\
            \frac{\partial f_m}{\partial x_1} (\ol{a}) & \ldots & \frac{\partial f_m}{\partial x_d} (\ol{a})
        \end{pmatrix}
    \]

    Если \( f \in D(\ol{a}) \), то \( J_f (\ol{a}) \) --- матрица \( df(\ol{a}) \)
\end{defn}

\newpage

\subsection{Достаточное условие дифференцируемости}

\begin{thrm}{}{}
    Пусть \( f : \RR^d \to \RR \)

    Пусть в некоторой окрестности \( U(\ol{a}) \)
    существуют \( \frac{\partial f}{\partial x_j} (\ol{x}) \ \forall j = 1, \ldots, d \),

    Причем \( \frac{\partial f}{\partial x_j} \in C(\ol{a}) \).

    Тогда \( f \in D(\ol{a}) \).
\end{thrm}

\begin{example}{}{}
    \[
        f(x, y) = \begin{cases}
            (x^2 + y^2) \sin \frac{1}{x^2 + y^2}, x^2 + y^2 \neq 0
            \\
            0, x = y = 0
        \end{cases}
    \]

    \( f \in D(\ol{0}) \), но не удовлетворяет условиям теоремы.
\end{example}

Приведем доказательство для \( d = 2 \):
\begin{gather*}
    f(a_1 + h_1, a_2 + h_2) - f(a_1, a_2)
    =
    \\
    =
    f(a_1 + h_1, a_2 + h_2) - f(a_1, a_2 + h_2) + f(a_1, a_2 + h_2) - f(a_1, a_2)
    =
    \\
    =
    f_1'(a_1 + \theta_1 h_1, a_2 + h_2) \cdot h_1 + f_2'(a_1, a_2 + \theta h_2) \cdot h_2
    =
    \\
    =
    f_1'(a_1, a_2) h_1 + f_2'(a_1, a_2) h_2
        + \\
        + ( f_1'(a_1 + \theta_1 h_1, a_2 + h_2) - f_1'(a_1, a_2) ) \cdot h_1
        + ( f_2'(a_1, a_2 + \theta h_2) - f_2'(a_1, a_2) ) \cdot h_2
    =
    \\
    =
    df(\ol{a}) \ol{h} + || \ol{h} || \alpha(\ol{h})
\end{gather*}

Для большего числа переменных нужно повторить рассуждение,
представив функцию как сумму изменений по каждой координате по очереди.

\subsection{Дифференцируемость сложного отображения и инвариантность формы первого дифференциала}

\begin{prop}{}{}
    Пусть \( L : \RR^d \to \RR^m \) --- линейное отображение.

    Тогда \( \exists C > 0 : \forall \ol{x} \in \RR^d \: || L \ol{x} || \leq C || \ol{x} || \)
\end{prop}

Пусть \( \ol{e}_1, \ldots, \ol{e}_d \) --- базис в \( \RR^d \),
а \( \ol{e}_1', \ldots, \ol{e}_m' \) --- базис в \( \RR^m \).
Пусть \( \ol{x} = \sum\limits_{j = 1}^{d} x_j \ol{e}_j \):
\[
    || L \ol{x} ||
    \leq
    \sum\limits_{j = 1}^{d} |x_j| \cdot || L \ol{e}_j ||
    \leq
    \max\limits_{j = 1, \ldots, d} |x_j| \cdot \sum\limits_{j = 1}^{d} || L \ol{e}_j ||
    \leq
    || \ol{x} || \cdot \sum\limits_{j = 1}^{d} || L \ol{e}_j ||
\]

Возьмем \( C = \sum\limits_{j = 1}^{d} || L \ol{e}_j || \) --- доказали.

\begin{thrm}{Дифференцируемость сложного отображения}{}
    Пусть \( f : \RR^d \to \RR^m \), \( f \in D(\ol{a}) \), \( g : \RR^m \to \RR^n \), \( g \in D(f(\ol{a})) \)

    Тогда \( g \circ f \in D(a) \), причем \( d (g \circ f) (\ol{a}) = dg(f(\ol{a})) \circ df (\ol{a}) \)
\end{thrm}

\begin{gather*}
    f(\ol{a} + \ol{h}) - f(\ol{a}) = df(\ol{a}) \ol{h} + || \ol{h} || \cdot \alpha (\ol{h})
    \\
    g(f(\ol{a}) + \ol{a}) - g(f(\ol{a})) = dg(f(\ol{a})) \ol{q} + || \ol{q} || \cdot \beta (\ol{q})
\end{gather*}

Доопределим \( \beta(\ol{0}) = 0 \).

Теперь подставим \( \ol{q} = f(\ol{a} + \ol{h}) - f(\ol{a}) = df(\ol{a}) \ol{h} + || \ol{h} || \cdot \alpha (\ol{h}) \):
\begin{gather*}
    g(f(\ol{a} + \ol{h})) - g(f(\ol{a}))
    =
    \\
    =
    dg(f(\ol{a})) ( df(\ol{a}) \ol{h} + || \ol{h} || \cdot \alpha (\ol{h}) )
        +
        || f(\ol{a} + \ol{h}) - f(\ol{a}) || \cdot \beta (f(\ol{a} + \ol{h}) - f(\ol{a}))
    =
    \\
    =
    (dg(f(\ol{a})) \circ df (\ol{a})) ( \ol{h} )
        +
        || \ol{h} || \cdot dg ( f(\ol{a}) ) ( \alpha(\ol{h}) )
        +
        \\
        +
        || f(\ol{a} + \ol{h}) - f(\ol{a}) || \cdot \beta (f(\ol{a} + \ol{h}) - f(\ol{a}))
    =
    \\
    =
    (dg(f(\ol{a})) \circ df (\ol{a})) ( \ol{h} ) + || \ol{h} || \cdot \gamma (\ol{h})
\end{gather*}

Проверим:
\[
    \gamma(\ol{h})
    =
    dg ( f(\ol{a}) ) ( \alpha(\ol{h}) )
    +
    \frac{|| f(\ol{a} + \ol{h}) - f(\ol{a}) ||}{|| \ol{h} ||} \cdot \beta (f(\ol{a} + \ol{h}) - f(\ol{a}))
\]

По предположению выше:
\[
    || dg ( f(\ol{a}) ) ( \alpha(\ol{h}) ) || \leq C || \alpha ( \ol{h} ) ||
\]

Значит это слагаемое стремится к нулю при \( \ol{h} \to \ol{0} \).

\begin{gather*}
    \frac{|| df(\ol{a}) \ol{h} + || \ol{h} || \cdot \alpha(h) ||}{|| \ol{h} ||} \cdot || \beta (f(\ol{a} + \ol{h}) - f(\ol{a})) ||
    \leq
    \\
    \leq
    \frac{C_1 || \ol{h} || + || \ol{h} || \cdot || \alpha(h) ||}{|| \ol{h} ||} \cdot || \beta (f(\ol{a} + \ol{h}) - f(\ol{a})) ||
    =
    \\
    =
    (C_1 + || \alpha(\ol{h}) ||) \cdot || \beta (f(\ol{a} + \ol{h}) - f(\ol{a})) ||
\end{gather*}

Получили константу на бесконечно малую.

А значит \( \gamma(\ol{h}) \to 0 \) при \( \ol{h} \to \ol{0} \)

\begin{coroll}{}{}
    \( J_g (f(\ol{a})) \cdot J_f (\ol{a}) \) --- матрица \( dg(f(\ol{a})) \circ df (\ol{a}) \)
\end{coroll}

Пусть \( f : \RR \to \RR^m \), \( g : \RR^m \to \RR \).

\( (g \circ f)' (a) = g_1'(f(a)) \cdot f_1'(a) + \ldots + g_m'(f(a)) \cdot f_m'(a) \),
где \( g_i' \) --- производная \( g \) по \( i \)-й координате, а \( f_i \) --- \( i \)-я координатная функция.

\begin{thrm}{Инвариантность формы первого дифференциала}{}
    Пусть \( f : \RR^d \to \RR^m \), \( f \in D(\ol{a}) \), \( g \RR^m \to \RR \), \( g \in D(\ol{y}) \), где \( \ol{y} = f(\ol{a}) \).

    Если в \( dg(\ol{y}) = \sum\limits_{k = 1}^{m} \frac{\partial g}{\partial y_k} (\ol{y}) dy_k \)
    вместо независимых переменных \( y_k \) подставить функции \( f_k(\ol{x}) \),
    то при \( \ol{x} = \ol{a} \) полученное выражение является \( d(g \circ f)(\ol{a}) \).
\end{thrm}

\begin{gather*}
    d(g \circ f)(\ol{a})
    =
    \sum\limits_{j = 1}^{d} (g \circ f)_j' (\ol{a}) dx_j
    =
    \sum\limits_{j = 1}^{d}
        \left(
            \sum\limits_{k = 1}^{m}
            \left[
                \frac{\partial g}{\partial y_k} (f(\ol{a})) \cdot \frac{\partial f_k}{\partial x_j} (\ol{a})
            \right]
        \right) dx_j
    =
    \\
    =
    \sum\limits_{k = 1}^{m}
        \left(
            \frac{\partial g}{\partial y_k} (f(\ol{a}))
            \cdot
            \sum\limits_{j = 1}^{d}
                 \frac{\partial f_k}{\partial x_j} (\ol{a})
                 dx_j
        \right)
    =
    \sum\limits_{k = 1}^{m} \frac{\partial g}{\partial y_k} (f(\ol{a})) df_k (\ol{a})
\end{gather*}

\subsection{Правила дифференцирования}

\begin{prop}{Правила дифференцирования}{}
    Пусть \( f, g : \RR^d \to \RR \), и \( f, g \in D(\ol{a}) \)

    \begin{itemize}
        \item
            \( (\alpha f + \beta g) \in D(\ol{a}) \)
        \item
            \( fg \in D(\ol{a}) \)
        \item
            При \( g(\ol{a}) \neq 0 \) \( \frac{f}{g} \in D(\ol{a}) \)
    \end{itemize}

    Сами правила:
    \begin{itemize}
        \item
            \( d (\alpha f + \beta g) ( \ol{a} ) = \alpha df (\ol{a}) + \beta dg (\ol{a}) \)
        \item
            \( d (fg) ( \ol{a} ) = f(\ol{a}) dg(\ol{a}) + g(\ol{a}) df (\ol{a}) \)
        \item
            \( d \left( \frac{f}{g} \right) ( \ol{a} ) = \frac{g(\ol{a}) df (\ol{a}) - f(\ol{a}) dg (\ol{a}) }{g^2(\ol{a})} \)
    \end{itemize}
\end{prop}

Проведем доказательство на примере произведения.

Пусть \( \varphi(x, y) = xy \Rightarrow d \varphi (x, y) = x dy + y dx \).

Если \( x = f(\ol{a}) \), \( y = g(\ol{a}) \),
то \( d (fg) (\ol{a}) = f (\ol{a}) dg (\ol{a}) + g (\ol{a}) df (\ol{a}) \).

\begin{example}{}{}
    \( g \in \RR^m \to \RR \), \( g \in D(\ol{b}) \).

    Пусть \( \ol{v} \in \RR^m \) и
    \(
        \ol{v} = \begin{pmatrix}
            \cos \alpha_1
            \\
            \vdots
            \\
            \cos \alpha_m
        \end{pmatrix}
    \), где \( \alpha_j \) --- угол между \( \ol{v} \) и осью \( y_j \).

    В таком случае \( || \ol{v} || = 1 \).

    \begin{gather*}
        g(\ol{b} + t \ol{v}) = g(b_1 + t v_1, \ldots, b_d + t v_m) = u(t)
        \\
        f(t) = \begin{pmatrix}
            b_1 + t \cos \alpha_1
            \\
            \vdots
            \\
            b_d + t \cos \alpha_m
        \end{pmatrix}
    \end{gather*}

    \begin{gather*}
        \frac{\partial g}{\partial \ol{v}} ( \ol{b} )
        =
        u' (0)
        =
        \frac{\partial g}{\partial y_1} (f(0)) \cdot (b_1 + t \cos \alpha_1)'
            +
            \ldots
            +
            \frac{\partial g}{\partial y_m} (f(0)) \cdot (b_m + t \cos \alpha_m)'
        \\
        \frac{\partial g}{\partial \ol{v}} ( \ol{b} )
        =
        \frac{\partial g}{\partial y_1} (f(0)) \cos \alpha_1
            +
            \ldots
            +
            \frac{\partial g}{\partial y_m} (f(0)) \cos \alpha_m
        =
        dg (\ol{b})(\ol{v})
        =
        (\nabla g(\ol{b}), \ol{v})
    \end{gather*}
\end{example}

\begin{example}{}{}
    \( g(y_1, y_2, \ldots, y_m) = g(f_1(x_1, \ldots, x_d), \ldots, f_m(x_1, \ldots, x_d)) \)

    \( g_{x_1}' = g_1' \cdot {f_1}_{x_1}' + \ldots + g_m' \cdot {f_m}_{x_1}' \)

    \( g : \RR^3 \to \RR \), \( f : \RR^2 \to \RR^3 \)

    \(
        f = \begin{pmatrix}
            x_1 + x_2
            \\
            x_1 x_2
            \\
            x_2 \sin x_1
        \end{pmatrix}
    \)
    \begin{gather*}
        d (g \circ f) (\ol{x})
        =
        g_{y_1}' dy_1 + g_{y_2}' dy_2 + g_{y_3} dy_3
        =
        \\
        =
        g_1' (dx_1 + dx_2) + g_2' (x_1 dx_2 + x_2 dx_1) + g_3' (x_2 \cos x_1 dx_1 + \sin x_1 dx_2)
        =
        \\
        =
        dx_1 (g_1' + x_2 g_2' + g_3' \cdot x_2 \cos x_1) + dx_2 (g_1' + x_1 g_2' + g_3' \sin x_1)
    \end{gather*}
\end{example}

\begin{thrm}{Дифференцируемость обратного отображения}{}
    Пусть \( f : \RR^d \to \RR^d \) биекция \( U(\ol{a}) \) на \( V(f(\ol{a})) \).

    Пусть \( f \in D(\ol{a}) \).

    Пусть \( df(\ol{a}) \) обратимо.

    Пусть \( \exists f^{-1} : V(f(\ol{a})) \to U(\ol{a}) \), причем \( f^{-1} \in C(\ol{a}) \).

    Тогда \( f^{-1} \in D(f(\ol{a})) \), причем \( df^{-1} (f(\ol{a})) = (df(\ol{a}))^{-1} \)
\end{thrm}

Если \( f^{-1} \in D(f(\ol{a})) \), то \( \forall \ol{q} \in \RR^d \) из некоторой проколотой окрестности нуля
выполняется следующее:
\[
    f^{-1} (f(\ol{a}) + \ol{q}) - f^{-1} (f(\ol{a})) = df^{-1} (\ol{a}) \ol{q} + \beta (\ol{q}) || \ol{q} ||
    \quad
    \beta(\ol{q}) \to \ol{0}, || \ol{q} || \to 0
\]

Покажем, что
\[
    \lim\limits_{|| \ol{q} || \to 0} \frac{|| f^{-1} (f(\ol{a}) + \ol{q}) - f^{-1} (f(a)) - df(\ol{a})^{-1} \ol{q} ||}{|| \ol{q} ||} \to 0
\]

Пусть \( \ol{h} = f^{-1} ( f(\ol{a}) + \ol{q} ) - f^{-1} ( f(\ol{a}) ) \).

Тогда \( \ol{q} = f(\ol{a} + \ol{h}) - f(\ol{a}) = df(\ol{a}) \ol{h} + || \ol{h} || \cdot \alpha ( \ol{h} ) \).

При достаточно малых \( || q || \) \( \ol{a} + \ol{h} \in U(\ol{a}) \), так как \( f^{-1} \in C(\ol{a}) \).

Аналогично \( f(\ol{a}) + \ol{q} \in V(f(\ol{a})) \)
\begin{gather*}
    \frac
    {|| \ol{h} - df(\ol{a})^{-1} (df(\ol{a}) \ol{h} + || \ol{h} || \alpha ( \ol{h} )) ||}
    {|| df(\ol{a}) \ol{h} + || \ol{h} || \cdot \alpha ( \ol{h} ) ||}
    =
    \\
    =
    \frac
    {|| \ol{h} || \cdot || df(\ol{a})^{-1} \alpha(\ol{h}) || }
    {|| df(\ol{a}) \ol{h} + || \ol{h} || \cdot \alpha ( \ol{h} ) ||}
\end{gather*}

Ранее выясняли, что \( \exists C > 0 : || (df (\ol{a})^{-1} \ol{x} ) || \leq C || \ol{x} || \).

Подставим \( \ol{x} = df (\ol{a}) \ol{h} \): \( || \ol{h} || \leq C || df(\ol{a}) \ol{h} || \).

Вернемся:
\begin{gather*}
    \frac
    {|| \ol{h} || \cdot || df(\ol{a})^{-1} \alpha(\ol{h}) ||}
    {|| df(\ol{a}) \ol{h} + || \ol{h} || \cdot \alpha ( \ol{h} ) ||}
    \leq
    \\
    \leq
    \frac
    {|| \ol{h} || \cdot C \cdot || \alpha(\ol{h}) ||}
    {|| \ol{h} || \cdot \left| || \frac{1}{C} \frac{\ol{h}}{|| \ol{h} ||} || - || \alpha(h) || \right| }
    =
    \frac
    {C \cdot || \alpha(\ol{h}) ||}
    {\left|\frac{1}{C} - || \alpha(h) || \right| }
\end{gather*}

Нетрудно видеть, что полученная дробь стремится к нулю при \( \ol{h} \to \ol{0} \).

\subsection{Теорема о неявной функции}

\begin{example}{}{}
    Пусть \( x^2 + y^2 - 1 = 0 \).

    Если \( x \) фиксированно, то решим уравнение относительно \( y \): \( y = \pm \sqrt{1 - x^2} \).
\end{example}

\begin{defn}{Неявная функция}{}
    Если дано уравнение \( F(x, y) = 0 \), из которого в окрестности некоторой точки \( (x_0, y_0) \)
    можно явно найти \( y \) и определить его единственным образом,
    то говорят, что \( y \) --- неявная функция от \( x \).
\end{defn}

\begin{thrm}{Теорема о неявной функции}{}
    Пусть
    \begin{enumerate}
        \item
            \( F : \RR^2 \to \RR \) непрерывна в окрестности \( \Omega \) точки \( (x_0, y_0) \).
        \item
            \( F(x_0, y_0) = 0 \)
        \item
            \( F_x', F_y' \in C(\Omega) \)
        \item
            \( F_y' (x_0, y_0) \neq 0 \)
    \end{enumerate}

    Тогда \( \exists! \) функция \( y = f(x) \),
    определенная в некоторой \( U_{\delta} (x_0) \):
    \begin{enumerate}
        \item
            \( f(x_0) = y_0 \)
        \item
            \( \forall x \in U_{\delta} (x_0) \: F(x, f(x)) = 0 \)
        \item
            \( f \in D(U_{\delta} (x_0)) \) и \( f'(x) = - \frac{F_x'(x, y) |_{y = f(x)} }{F_y'(x, y) |_{y = f(x)} } \)
    \end{enumerate}
\end{thrm}

Считаем, что \( F_y' (x_0, y_0) > 0 \), так как иначе можно перейти к \( -F \).

Так как \( F_y' \in C(\Omega) \) и \( F_y'(x_0, y_0) > 0 \),
то существуют \( a > 0 \) и \( b > 0 \) такие, что \( F_y' > m > 0 \)
в прямоугольнике \( p = [x_0 - a, x_0 + a] \times [y_0 - b, y_0 + b] \).

Так как \( F_y' > 0 \) на \( p \), \( F(x_0, y) \uparrow \) на \( [y_0 - b, y_0 + b] \).

Так как \( F(x_0, y_0) = 0 \), \( F(x_0, y_0 - b) < 0 \), \( F(x_0, y_0 + b) > 0 \).

Можно выбрать достаточно малое \( a \) такое, что два вышеописанных утверждения
будут верны для \( x \in [x_0 - a, x_0 + a] \).
А значит на каждой ``вертикали'' найдется ровно одна точка \( y = f(x) \), где \( F \) обращается в ноль.

Для такой \( f(x) \) верно, что \( f(x_0) = y_0 \).
Единственна же она, так как для каждого фиксированного \( x \) нужный \( y \) единственен.

Возьмем \( \delta > 0 : (x_0 - \delta, x_0 + \delta) \subset [x_0 - a, x_0 + a] \).
Этот интервал и есть \( U_{\delta} (x_0) \).

Пусть \( x, c \in U_{\delta} (x_0) \) и \( y_c = f(c), y = f(x) \).

\( F(c, y_c) = F(x, y) = 0 \).

Для функции \( h(t) = F(c + t(x - c), y_c + t(y - y_c)) \) имеем:
\( h(0) = h(1) = 0 \) и \( h \in C[0, 1] \), \( h \in D(0, 1) \).

\( h'(t) = F_1' \cdot (x - c) + F_2' \cdot (y - y_c) \).

По теореме Ролля \( \exists \xi \in (0, 1) : h'(\xi) = 0 \)

\begin{gather*}
    0 = F_1' \cdot (x - c) + F_2' (y - y_c)
    \\
    \Updownarrow
    \\
    \frac{y - y_c}{x - c} = - \frac{F_1'(c + \xi (x - c), y_c + \xi (y - y_c))}{F_2'(c + \xi (x - c), y_c + \xi (y - y_c))}
\end{gather*}

Так как \( F_1', F_2' \) принадлежат \( C(p) \), \( \exists M = \max\limits_{p} | F_1 (x, y) | \).
Еще \( F_y'(x, y) \geq m > 0 \).

Откуда \( \left| \frac{y - y_c}{x - c} \right| \leq \frac{M}{m} \).

\( y - y_c = (x - c) \frac{y - y_c}{x - c} \Rightarrow \) при \( x \to c \) выполнено \( y \to y_c \).
То есть \( \lim\limits_{x \to c} f(x) = f(c) \), то есть функция непрерывна в любой точке окрестности.
Если \( x \to c \), то множитель при \( \xi \) стремится к нулю, поэтому
\[
    \lim\limits_{x \to c} \frac{f(x) - f(c)}{x - c} = - \frac{F_x'(c, f(c))}{F_y'(c, f(c))}
\]

Можем взять в качестве \( c \) любую точку интервала \( \Rightarrow \) доказали теорему.

\begin{example}{}{}
    Вернемся к \( F(x, y) = x^2 + y^2 - 1 \)

    \( F_y' = 2y \neq 0 \ \forall (x, y) : y \neq 0 \)

    \( f'(x) = -\frac{F'(x, y)}{F_y'(x, y)} = -\frac{2x}{2f(x)} = -\frac{x}{f(x)} \)

    \( f(x) = \pm \sqrt{1 - x^2} \)
\end{example}

Пусть \( (\ol{x}, y) = (x_1, \ldots, x_{d - 1}, y) \in \RR^d \)

\begin{thrm}{Общая теорема о неявной функции}{}
    Пусть
    \begin{enumerate}
        \item
            \( F : \RR^d \to \RR \) непрерывна в окрестности \( \Omega \) точки \( (\ol{x}_0, y_0) \).
        \item
            \( F(\ol{x}_0, y_0) = 0 \)
        \item
            \( F_{x_1}', \ldots, F_{x_{d - 1}}', F_y' \in C(\Omega) \)
        \item
            \( F_y' (\ol{x}_0, y_0) \neq 0 \)
    \end{enumerate}

    Тогда \( \exists! \) функция \( y = f(\ol{x}) \),
    определенная в некоторой \( U_{\delta} (\ol{x}_0) \):
    \begin{enumerate}
        \item
            \( f(\ol{x}_0) = y_0 \)
        \item
            \( \forall \ol{x} \in U_{\delta} (\ol{x}_0) \: F(\ol{x}, f(\ol{x})) = 0 \)
        \item
            \( f \) имеет все частные производные по \( x_1, \ldots, x_{d - 1} \) в \( U_{\delta}(\ol{x}_0) \)

            Причем \( f_{x_j}' (\ol{x}) = -\frac{F_{x_j}' (\ol{x}, f(\ol{x}))}{F_y' (\ol{x}, f(\ol{x}))} \)
    \end{enumerate}
\end{thrm}

Идея доказательства аналогична.
